
\subsubsection{3.1 Overview}

The goal is to detect when classifier commitment to spurious features accelerates. We hypothesize this manifests as a localized acceleration in worst-group accuracy decline---a negative second derivative---rather than gradual monotonic decline.

The detection pipeline consists of three stages: (1) dense WGA tracking during training, (2) smoothing to reduce noise, and (3) second derivative computation with peak detection.

\subsubsection{3.2 Worst-Group Accuracy Tracking}

Worst-group accuracy is computed at every epoch during training. For a dataset with group annotations (spurious attribute $\times$ label), WGA is the minimum accuracy across all groups:

$$\text{WGA} = \min_{g \in \mathcal{G}} \frac{1}{|D_g|} \sum_{(x,y) \in D_g} \mathbf{1}[\hat{y}(x) = y]$$

Dense checkpointing (every epoch) is required to avoid missing the crystallization peak.

\subsubsection{3.3 Smoothing}

Raw WGA curves exhibit epoch-to-epoch noise due to batch sampling variance. A rolling average with a window of 5 epochs is applied. Window sizes of 3, 5, and 7 epochs all achieve detection, with 5 epochs providing the highest signal-to-noise ratio in our experiments.

\subsubsection{3.4 Second Derivative Detection}

The second derivative is computed using central differences on the smoothed WGA. A significant negative peak in $d^2WGA/dt^2$ indicates accelerating decline. Detection parameters: prominence threshold >0.005, search within first 50\% of training, SNR threshold >2.0.

\subsubsection{3.5 Linear Probe Analysis}

To verify whether crystallization represents classifier commitment rather than representation degradation, linear probes are trained on frozen penultimate layer activations to predict both spurious and core attributes.

\subsubsection{3.6 Gradient Ratio Tracking}

To verify the causal role of gradient starvation, the gradient ratio between minority and majority group examples is tracked. If gradient starvation causes crystallization, the gradient ratio inflection should precede the WGA crystallization peak.

